\scanpage{359}
\mspage{98}

Alors le foncteur $(\unsure{\alpha|*})$ (p. 94) \deleted{\unsure{peut}} se factorise par la sous-catégorie $\underline{\mathrm{Hom}}_{\mathrm{Cat}/X\times X}(\,.\,,\,.\,)$ de $\underline{\mathrm{Hom}}_{\mathrm{Cat}}(\,.\,,\,.\,)$, formée des foncteurs qui commutent aux foncteurs de projection dans $X\times X$, et les hom. entre tels foncteurs qui s'effectuent par des flèches ``verticales'' \unsure{de|dans} la catégorie but (i.e. par des flèches au dessus des flèches identiques de la catégorie base $X\times X$).
\[
\underline{\mathrm{Hom}}(\tau,\tau')\longrightarrow\underline{\mathrm{Hom}}_{\mathrm{Cat}/X\times X}
\bigl(\underline{\mathrm{Ch}}_{\Sigma,\tau'}(X),\underline{\mathrm{Ch}}_{\Sigma,\tau}(X)\bigr)
\]
\unsure{:}

Donc dans l'énoncé du théorème d'homotopie, on peut préciser que l'homotopie se fait dans cette sous-catégorie, elle s'effectue par \deleted{\unsure{une|des}}\unsure{des|une} \unsure{suites|suite|bouts} d'homotopies simples qui sont des morphismes de foncteurs entre \deleted{$X\times X$-}catégories au dessus de $X\times X$. Cela implique alors l'énoncé d'homotopie sur les $\varphi_{X,a,b}^*$.

% archive sheet number: 358
